% ==============================================================================
% Inside Synthetic Dollar Liquidity and the Microstructure of FX Swap Markets
% Formal Theoretical Framework — Version 2
% Revised by: Theorist agent, 2026-05-18
%
% REVISION LOG — Round 2 (2026-05-18)
% Five blocking fixes applied per Theorist-Critic report (score: 24/100):
%   Fix 1: Prop 1 Step 2 — added quality-adjusted price $\tilde{R}_{I,t}$
%           and IFT-based derivation of $\partial\mu_t^*/\partial\chi_t < 0$.
%           Note propagated to Props 2, 4, 5.
%   Fix 2: Prop 7(b) — removed Poincaré-Hopf; corrected boundary-condition
%           sign; replaced with correct IVT argument on $F(\mu)=\Psi(\mu)-\mu$.
%   Fix 3: Prop 3 — added $\Phi''>0$ as explicit sufficient condition in
%           statement; added Remark on when it holds and when it may fail.
%   Fix 4: Prop 0 (Existence) — added global existence paragraph via IVT on
%           the aggregate net-demand affine function.
%   Fix 5: Prop 6 — downgraded $|\beta_{FB}|<\beta_{HF}$ from derived result
%           to Maintained Hypothesis with explicit empirical motivation.
%
% REVISION LOG — Round 3 (2026-05-18)
% Three targeted fixes applied per Theorist-Critic Round 2 report:
%   Fix 1 (CRITICAL): Prop 1 Step 2 — removed erroneous derivation via the
%           price-recovery identity $\mu_t = \chi_t(\tilde\mu_t^*+R_{O,t})-R_{O,t}$
%           (which gave $\partial\mu_t^*/\partial\chi_t = \tilde{R}_{I,t}^* > 0$,
%           contradicting the IFT result). Removed the contradicting equation
%           eq:dmu-dchi and the "Wait —" acknowledgment entirely. Added a
%           clarifying note explaining why the identity does not yield the
%           comparative static ($\tilde\mu_t^*$ is defined at fixed $\chi_t$).
%           Only the IFT-based argument (giving $\partial\mu_t^*/\partial\chi_t
%           < 0$) is retained.
%   Fix 2 (MAJOR): Prop 7(b) — replaced informal local-slope argument for
%           existence of second fixed point (which incorrectly inferred
%           $F(\hat\mu) < 0$ from $F'(\hat\mu) < -1$ alone) with a rigorous
%           global condition: $|\Lambda(\mu)| > 1$ for all $\mu \in (c^{HF},
%           \mu_1^*)$ ensures $F$ descends by more than its initial height,
%           crossing zero before $\mu_1^*$. Added Remark rem:global-lambda
%           on sufficiency vs.\ necessity of the global condition.
%   Fix 3 (MAJOR): Added Remark rem:B-negative (Maintained Hypothesis:
%           Downward-Sloping Aggregate Demand) immediately after
%           Assumption ass:clearing, introducing $\mathcal{B} \equiv \sum_s
%           \beta_s < 0$ as the explicit calligraphic symbol for the aggregate
%           slope (distinct from $B_t$ capacity). Updated all references in
%           the Prop 0 proof, Definition def:price-eq, Remark rem:price-impact,
%           Prop 1 proof, and Remark on convexity conditions from plain $B$ to
%           $\mathcal{B}$.
% ==============================================================================

\documentclass[12pt]{article}

% ====== Page Layout ======
\usepackage[left=1.0in,right=1.0in,top=1.0in,bottom=1.0in]{geometry}
\usepackage{setspace}
\doublespacing
\usepackage{fancyhdr}
\pagestyle{fancy}
\fancyhf{}
\fancyfoot[C]{\thepage}
\renewcommand{\headrulewidth}{0pt}

% ====== Typography ======
\usepackage{lmodern}
\usepackage{microtype}
\usepackage[T1]{fontenc}

% ====== Math ======
\usepackage{amssymb,amsmath,amsfonts,mathtools}
\usepackage{dsfont}
\usepackage{amsthm}

\newtheorem{theorem}{Theorem}
\newtheorem{proposition}{Proposition}
\newtheorem{corollary}{Corollary}[proposition]
\newtheorem{lemma}{Lemma}
\newtheorem{definition}{Definition}
\newtheorem{assumption}{Assumption}
\newtheorem{remark}{Remark}

\theoremstyle{definition}
\newtheorem{example}{Example}

\DeclareMathOperator*{\argmax}{arg\,max}
\DeclareMathOperator*{\argmin}{arg\,min}
\newcommand{\norm}[1]{\left\lVert #1 \right\rVert}
\newcommand{\1}[1]{\mathds{1}\left[#1\right]}
\newcommand{\E}{\mathbb{E}}
\newcommand{\R}{\mathbb{R}}

% ====== Tables ======
\usepackage{array,booktabs,makecell}
\usepackage[flushleft]{threeparttable}
\usepackage{tabularray}
\UseTblrLibrary{booktabs}

% ====== Figures ======
\usepackage{graphicx}
\usepackage{caption}
\captionsetup{font=small,labelfont=bf,justification=justified}

% ====== Lists ======
\usepackage{enumitem}

% ====== Hyperref (second-to-last) ======
\usepackage[hidelinks,breaklinks,colorlinks=true,
            linkcolor=black,citecolor=black,urlcolor=black]{hyperref}

% ====== Cleveref (after hyperref) ======
\usepackage[nameinlink]{cleveref}

% ====== Section headings ======
\usepackage{titlesec}

% ==============================================================================
\title{Inside Synthetic Dollar Liquidity and the Microstructure\\
       of FX Swap Markets: A Formal Theoretical Framework%
       \thanks{This document contains the formal theory section for the paper
       ``Inside Synthetic Dollar Liquidity and the Microstructure of FX Swap
       Markets.'' Notation is consistent with the empirical roadmap v8.}}
\author{}
\date{May 2026}
% ==============================================================================

\begin{document}

\maketitle
\thispagestyle{empty}
\newpage
\setcounter{page}{1}

% ==============================================================================
\section{Introduction and Overview}
\label{sec:intro}
% ==============================================================================

This document formalizes the theoretical framework underlying the paper. The
central claim is:

\begin{quote}
  Dollar funding crises are not runs against the dollar. They are runs \emph{from}
  inside synthetic dollar liquidity into outside safe dollar liquidity. The FX
  swap market is where this conversion is priced and rationed.
\end{quote}

The framework combines three building blocks: (i) a CES aggregator that makes
inside and outside dollar liquidity imperfect substitutes, with the degree of
substitutability governed by an endogenous quality parameter; (ii) a
search-and-matching market for conversion intermediaries, in which market
tightness determines conversion probability and hence quality; and (iii) an
equilibrium demand system across institutional sectors. The main results are
seven propositions about the equilibrium funding spread $\mu_t$, its
comparative statics in capacity and demand, its non-linear behavior at high
market tightness, its seasonal amplification at regulatory reporting dates, the
stabilizing role of central-bank swap lines, the heterogeneous supply behavior
of intermediary types, and the conditions under which a self-reinforcing run
equilibrium exists.

\bigskip

\noindent\textbf{Relationship to the empirical design.} Each proposition maps to
one or more empirical predictions in the v8 design (labeled P1--P9 there).
\Cref{tab:theory-empirics} at the end of this document makes this mapping
explicit.

% ==============================================================================
\section{Primitives and Notation}
\label{sec:notation}
% ==============================================================================

\subsection{Economic environment}

Time is discrete, indexed by $t \in \{1, 2, \ldots\}$. The economy is
populated by a unit mass of agents who require dollar liquidity to meet payment
obligations. Agents are heterogeneous in their institutional access to safe
dollar liquidity; we abstract from individual heterogeneity and characterize the
representative demand structure through a CES aggregator and a sectoral demand
system.

\paragraph{Endogenous variables.} At each date $t$ the following objects are
determined in equilibrium:
\begin{itemize}[noitemsep]
  \item $I_t \geq 0$: quantity of inside (synthetic) dollar liquidity demanded
  \item $O_t \geq 0$: quantity of outside (safe) dollar liquidity demanded
  \item $L_t \geq 0$: effective aggregate dollar liquidity service
  \item $\chi_t \in (0,1]$: quality of inside dollar liquidity (formerly $q_t$ in the CES layer; renamed to avoid conflict with sectoral quantities; see \cref{rem:notation})
  \item $p_t \in (0,1]$: matching probability for a seeker in the conversion market
  \item $\theta_t \geq 0$: market tightness
  \item $\mu_t \in \R$: equilibrium funding spread (cross-currency basis)
  \item $R_{I,t}, R_{O,t} \geq 0$: gross costs of inside and outside dollar liquidity
\end{itemize}

\paragraph{Exogenous variables and parameters.} At each date $t$:
\begin{itemize}[noitemsep]
  \item $U_t \geq 0$: mass of agents actively seeking to convert inside to outside
        dollar liquidity (demand-side market participant mass)
  \item $B_t \geq 0$: aggregate balance-sheet capacity of conversion intermediaries
  \item $\bar{L}_t > 0$: minimum liquidity requirement (exogenous, reflecting
        payment obligations)
  \item $\alpha_t \in (0,1)$: preference weight on outside safe dollar liquidity
  \item $\rho < 1$, $\rho \neq 0$: CES substitution parameter; $\sigma \equiv 1/(1-\rho) > 1$
  \item $\eta > 0$, $\gamma \in (0,1)$: matching technology parameters
\end{itemize}

\begin{remark}[Notation convention for quality vs.\ quantity]
  \label{rem:notation}
  The draft uses $q_t$ in two distinct senses: (i) as the quality parameter
  in the CES aggregator, and (ii) as net-demand quantities in the sectoral
  demand system (equation~\eqref{eq:net-demand}). These are conceptually
  distinct objects. This revision uses $\chi_t \in (0,1]$ exclusively for
  quality, and $\mathfrak{q}_{s,m,t}$ for sectoral net quantities in the
  demand system. The empirical counterpart of $\chi_t$ is price dispersion
  across sectors; the empirical counterpart of $\mathfrak{q}_{s,m,t}$ is
  \texttt{Q\_{s,m,t}} in v8. The notation \texttt{PriceUSD}\,$_{m,t}$ in the
  v8 corresponds to $\mu_t$ (or $\mu_{m,t}$ when indexed by currency pair $m$)
  throughout.
\end{remark}

% ==============================================================================
\section{Definitions}
\label{sec:definitions}
% ==============================================================================

\begin{definition}[Outside Safe Dollar Liquidity]
  \label{def:outside}
  \emph{Outside safe dollar liquidity}, denoted $O_t$, is a claim on dollars
  connected directly to the US institutional core: Federal Reserve reserve
  balances, highly liquid Treasury securities, domestic bank liabilities backed
  by FDIC insurance or discount window access, and dollar liquidity extended via
  central-bank swap lines. $O_t$ is measured in units of effective liquidity
  service at date $t$.
\end{definition}

\begin{definition}[Inside Synthetic Dollar Liquidity]
  \label{def:inside}
  \emph{Inside synthetic dollar liquidity}, denoted $I_t$, is dollar liquidity
  produced outside the US domestic banking system, primarily via FX swaps, offshore
  repo, Eurodollar deposits, and the wholesale funding operations of global
  banks. A unit of $I_t$ is obtained when a non-US entity delivers its domestic
  currency on the near leg of an FX swap and commits to return dollars on the
  far leg. $I_t$ is measured in gross notional USD units.
\end{definition}

\begin{definition}[Synthetic Dollar Quality]
  \label{def:quality}
  The \emph{synthetic dollar quality} $\chi_t \in (0,1]$ is the number of units
  of outside safe dollar liquidity service that one unit of $I_t$ delivers,
  given the current state of the conversion market. Formally,
  \[
    \chi_t = \chi(p_t), \quad \chi\colon [0,1] \to (0,1],
  \]
  where $\chi$ is continuously differentiable with $\chi'(p) > 0$ and
  $\chi(1) = 1$. When $p_t = 1$, every seeker finds a counterparty and inside
  dollar liquidity is a perfect substitute for outside safe dollar liquidity;
  when $p_t \to 0$, synthetic liquidity delivers arbitrarily little service.
\end{definition}

\begin{definition}[Effective Dollar Liquidity Service]
  \label{def:CES}
  The \emph{effective aggregate dollar liquidity service} is produced by a CES
  technology:
  \begin{equation}
    L_t = \bigl[\alpha_t O_t^{\rho} + (1-\alpha_t)(\chi_t I_t)^{\rho}\bigr]^{1/\rho},
    \label{eq:CES}
  \end{equation}
  where $\rho < 1$, $\rho \neq 0$, and $\sigma \equiv 1/(1-\rho) \in (1, \infty)$
  is the elasticity of substitution. The specification $\sigma > 1$ captures the
  normal-times approximate equivalence between $I_t$ and $O_t$ (they are gross
  substitutes) while allowing the effective marginal rate of substitution to
  diverge in stress when $\chi_t \ll 1$.
\end{definition}

\begin{definition}[Conversion Market and Matching Technology]
  \label{def:matching}
  The \emph{conversion market} is a search-and-matching market in which agents
  holding $I_t$ seek counterparties who can supply $O_t$. The aggregate matching
  function is
  \begin{equation}
    M_t = \eta\, U_t^{\gamma}\, B_t^{1-\gamma},
    \label{eq:matching}
  \end{equation}
  with $\eta > 0$, $\gamma \in (0,1)$. The matching function satisfies constant
  returns to scale in $(U_t, B_t)$. The \emph{matching probability} for a seeker
  is
  \begin{equation}
    p_t = \frac{M_t}{U_t} = \eta\left(\frac{B_t}{U_t}\right)^{1-\gamma}
        = \eta\,\theta_t^{-(1-\gamma)},
    \label{eq:match-prob}
  \end{equation}
  where \emph{market tightness} is $\theta_t \equiv U_t/B_t \geq 0$.
\end{definition}

\begin{definition}[Funding Spread]
  \label{def:spread}
  The \emph{funding spread} (or \emph{synthetic dollar premium}) is
  \[
    \mu_t \equiv R_{I,t} - R_{O,t},
  \]
  where $R_{I,t}$ and $R_{O,t}$ are the gross dollar costs of inside and outside
  dollar liquidity, respectively. When $\mu_t > 0$, inside dollar liquidity is
  more expensive than outside: agents pay a premium for synthetic dollars.
  Empirically, $\mu_{m,t} = \text{PriceUSD}_{m,t} = -\text{TIB}_{m,t}$ in
  currency pair $m$.
\end{definition}

\begin{definition}[Intermediary Capacity]
  \label{def:capacity}
  \emph{Intermediary capacity} is
  \begin{equation}
    B_t = B_t^{FBO} + B_t^{USbank} + B_t^{dealer} + B_t^{HF} + B_t^{CBswap},
    \label{eq:capacity}
  \end{equation}
  where each component is the balance-sheet capacity of the respective
  intermediary type (foreign bank office branches, US commercial banks, primary
  dealers, hedge funds, and central-bank swap-line facilities, respectively).
  Each component is non-negative and taken as exogenous within a period, though
  it responds to state variables across periods under extensions
  (\cref{sec:extensions}).
\end{definition}

% ==============================================================================
\section{Assumptions}
\label{sec:assumptions}
% ==============================================================================

\begin{assumption}[Liquidity Requirement]
  \label{ass:liq-req}
  At every date $t$, the representative agent must meet a minimum liquidity
  requirement $\bar{L}_t > 0$, which is exogenous and strictly positive.
  \emph{Interpretation:} $\bar{L}_t$ reflects inelastic payment obligations
  (settlement, margin calls, mandatory FX hedging ratios). It rules out the
  degenerate case in which agents can simply reduce their dollar footprint when
  conversion is costly; they cannot.
\end{assumption}

\begin{assumption}[Cost Minimization]
  \label{ass:cost-min}
  The representative agent takes $R_{I,t}$, $R_{O,t}$, and $\chi_t$ as given
  and solves:
  \begin{equation}
    \min_{O_t,\, I_t \geq 0} \quad R_{O,t}\, O_t + R_{I,t}\, I_t
    \quad \text{subject to} \quad L_t \geq \bar{L}_t.
    \label{eq:cost-min}
  \end{equation}
  \emph{Interpretation:} Agents are price-takers in both liquidity markets.
  This rules out strategic behavior by the representative agent; it is
  appropriate when any single institution is small relative to the market.
\end{assumption}

\begin{assumption}[Imperfect Substitution]
  \label{ass:sigma}
  The elasticity of substitution satisfies $\sigma \in (1, \infty)$, i.e.\
  $\rho \in (0,1)$. Both $O_t$ and $I_t$ are used in positive quantity in any
  interior solution.
  \emph{Interpretation:} $\sigma < \infty$ means agents cannot costlessly
  replace outside with inside liquidity, even at unchanged prices. $\sigma > 1$
  means they are gross substitutes (not complements), so price spreads induce
  substitution.
\end{assumption}

\begin{assumption}[Quality Function]
  \label{ass:quality}
  The quality function $\chi\colon[0,1]\to(0,1]$ is continuously
  differentiable, strictly increasing ($\chi'(p) > 0$ for all $p \in (0,1]$),
  and satisfies $\chi(1) = 1$ and $\lim_{p\to 0^+} \chi(p) = \chi_{\min} > 0$.
  \emph{Interpretation:} A lower bound $\chi_{\min} > 0$ prevents the model
  from generating a complete collapse of inside liquidity, which would be
  inconsistent with observed markets (even in March 2020, FX swaps continued to
  trade at positive volume). The assumption $\chi(1)=1$ normalizes the
  frictionless benchmark.
\end{assumption}

\begin{assumption}[Matching Function Properties]
  \label{ass:matching}
  The matching function $M_t = \eta U_t^{\gamma} B_t^{1-\gamma}$ satisfies:
  (i) $M_t \leq \min(U_t, B_t)$; (ii) $M_t$ is increasing and concave in each
  argument; (iii) $p_t = M_t/U_t$ is decreasing in $\theta_t = U_t/B_t$.
  \emph{Interpretation:} Condition (i) is a feasibility constraint: matches
  cannot exceed either side of the market. Condition (iii) is the central search
  externality: each additional seeker reduces the matching probability for all
  others.
\end{assumption}

\begin{assumption}[Spread Determination]
  \label{ass:spread}
  In equilibrium, the funding spread satisfies the indifference condition of
  the marginal agent: for any interior solution with $I_t > 0$ and $O_t > 0$,
  the ratio of marginal costs equals the ratio of marginal products:
  \begin{equation}
    \frac{R_{I,t}}{R_{O,t}} = \frac{\alpha_t}{\,1-\alpha_t\,} \cdot
    \left(\frac{\chi_t I_t}{O_t}\right)^{\rho-1} \cdot \chi_t^{\rho}.
    \label{eq:foc}
  \end{equation}
  \emph{Interpretation:} This is the first-order condition of
  \eqref{eq:cost-min}. At the optimum, the marginal cost of obtaining one unit
  of effective liquidity service via inside and outside channels must be equal.
\end{assumption}

\begin{assumption}[Market Clearing]
  \label{ass:clearing}
  For each currency pair $m$ and date $t$:
  \begin{equation}
    \sum_{s} \mathfrak{q}_{s,m,t} = 0,
    \label{eq:clearing}
  \end{equation}
  where $\mathfrak{q}_{s,m,t}$ is the signed net USD-taking of sector $s$ in
  market $m$ at date $t$, with $\mathfrak{q}_{s,m,t} > 0$ indicating net USD
  receipt (demand side) and $\mathfrak{q}_{s,m,t} < 0$ indicating net USD
  supply. \emph{Interpretation:} The FX swap market clears bilaterally: every
  dollar synthetic received on the near leg by one counterparty is supplied by
  another.
\end{assumption}

\begin{remark}[Maintained Hypothesis: Downward-Sloping Aggregate Demand]
  \label{rem:B-negative}
  We maintain that $\mathcal{B} \equiv \sum_s \beta_s < 0$: aggregate net
  demand for synthetic dollar liquidity is downward-sloping in the funding
  spread $\mu_{m,t}$. This condition holds if and only if the dollar synthetic
  liquidity market is a normal goods market --- higher cost reduces net demand.
  It is equivalent to $\sigma > 0$ in the CES aggregator and is testable: if
  the estimated $\hat{\mathcal{B}} \geq 0$ from the demand system, downward-sloping
  aggregate demand is rejected by the data and the equilibrium existence result
  fails. We treat $\mathcal{B} < 0$ as a maintained hypothesis verified ex post
  by the sign of $\hat{\mathcal{B}}$.
\end{remark}

\begin{assumption}[Intermediary Supply Behavior]
  \label{ass:supply}
  Each intermediary type's capacity responds to state variables as follows.
  (a) FBO branches:
  \begin{equation}
    B_t^{FBO} = \bar{B}^{FBO} - \psi S_t - \omega \mathds{1}[Q_t = 1] + \nu A_t^{reg},
    \label{eq:FBO-capacity}
  \end{equation}
  where $S_t \geq 0$ is a global financial stress index, $Q_t \in \{0,1\}$ is
  a quarter-end indicator, and $A_t^{reg} \geq 0$ is a relative regulatory
  advantage of FBO branches. All parameters $\psi, \omega, \nu \geq 0$.
  (b) Hedge funds:
  \begin{equation}
    B_t^{HF} = \phi\,\max(\mu_t - c^{HF}, 0) \cdot \kappa_t,
    \label{eq:HF-capacity}
  \end{equation}
  where $c^{HF} \geq 0$ is the hedge fund's entry cost for the basis trade and
  $\kappa_t \geq 0$ is their risk capacity (exogenous within a period).
  (c) Central-bank swap lines: $B_t^{CBswap}$ is set by policy and treated as
  an exogenous shifter in the static equilibrium; see \cref{sec:extensions} for
  the optimal intervention extension.
  \emph{Interpretation:} Assumption~\ref{ass:supply} makes explicit which
  forces compress capacity at quarter-end and in stress (FBO branches) and which
  forces deliver pro-cyclical supply (hedge funds, who earn more when the basis
  is wide but retreat when risk capacity is exhausted).
\end{assumption}

\begin{assumption}[Regularity for Comparative Statics]
  \label{ass:regularity}
  The equilibrium values of $(\mu_t, \chi_t, p_t, \theta_t)$ are continuously
  differentiable functions of $(U_t, B_t, \alpha_t, \bar{L}_t)$ in a neighborhood
  of any interior equilibrium. This holds when $\chi$ is $C^2$ and $B_t > 0$.
\end{assumption}

% ==============================================================================
\section{Equilibrium}
\label{sec:equilibrium}
% ==============================================================================

\begin{definition}[Static Equilibrium]
  \label{def:equilibrium}
  A \emph{static equilibrium} at date $t$ is a tuple
  \[
    \mathcal{E}_t = \bigl(O_t^*, I_t^*, L_t^*, \chi_t^*, p_t^*, \theta_t^*,
    \mu_t^*, R_{I,t}^*, R_{O,t}^*\bigr)
  \]
  such that:
  \begin{enumerate}[label=(\roman*)]
    \item \textbf{Cost minimization:} $(O_t^*, I_t^*)$ solve
          \eqref{eq:cost-min} given $(R_{O,t}^*, R_{I,t}^*, \chi_t^*)$, with
          $L_t^* = \bar{L}_t$ (the constraint binds by cost minimization with
          positive costs).
    \item \textbf{Matching:} $p_t^* = \eta(B_t/U_t)^{1-\gamma}$ and
          $\theta_t^* = U_t/B_t$.
    \item \textbf{Quality:} $\chi_t^* = \chi(p_t^*)$.
    \item \textbf{Spread:} $\mu_t^* = R_{I,t}^* - R_{O,t}^*$ satisfies the
          first-order condition~\eqref{eq:foc} at $(O_t^*, I_t^*, \chi_t^*)$.
    \item \textbf{Market clearing:} \eqref{eq:clearing} holds for all $m$.
    \item \textbf{Sectoral demand:} Each sector $s$ optimizes over signed
          positions $\mathfrak{q}_{s,m,t}$ given $\mu_{m,t}^*$ and sector-specific
          characteristics; see \cref{sec:demand-system}.
  \end{enumerate}
\end{definition}

\begin{proposition}[Equilibrium Existence]
  \label{prop:existence}
  Under Assumptions~\ref{ass:liq-req}--\ref{ass:clearing}, a static equilibrium
  $\mathcal{E}_t$ exists for any $(U_t, B_t) \in \R_{++}^2$ and any exogenous
  parameters $(\bar{L}_t, \alpha_t, \rho, \eta, \gamma) \in \R_{++} \times
  (0,1) \times (-\infty,1)\setminus\{0\} \times \R_{++} \times (0,1)$.
\end{proposition}

\begin{proof}
  We construct the equilibrium directly in three steps.

  \textit{Step 1 (Matching and quality).} Given $(U_t, B_t) \in \R_{++}^2$,
  define $\theta_t = U_t/B_t > 0$. By \eqref{eq:match-prob},
  $p_t = \eta\,\theta_t^{-(1-\gamma)} \in (0,\infty)$. Because $U_t, B_t > 0$
  and $\gamma \in (0,1)$, we have $p_t > 0$. If $p_t > 1$ (possible when
  $B_t \gg U_t$), we cap at $p_t = 1$ by defining $p_t = \min(\eta\,
  \theta_t^{-(1-\gamma)}, 1)$. By Assumption~\ref{ass:quality}, $\chi_t =
  \chi(p_t)$ is well-defined and in $(0,1]$.

  \textit{Step 2 (Cost minimization).} Given $\chi_t$ and $(R_{O,t}, R_{I,t})$
  with $R_{O,t}, R_{I,t} > 0$, the cost-minimization problem~\eqref{eq:cost-min}
  has a unique interior solution by strict convexity of the CES cost function
  (which holds because $\sigma > 1$ and $\chi_t > 0$). The unique solution
  satisfies the FOC~\eqref{eq:foc} and the constraint $L_t = \bar{L}_t$,
  yielding closed-form demands:
  \begin{align}
    O_t^* &= \bar{L}_t \cdot \frac{\alpha_t^{\sigma}}{
      \bigl[\alpha_t^{\sigma} R_{O,t}^{1-\sigma} +
      (1-\alpha_t)^{\sigma}(\chi_t R_{I,t})^{1-\sigma}\bigr]^{1/(1-\sigma)}\,
      R_{O,t}^{-\sigma}}, \label{eq:demand-O} \\
    I_t^* &= \bar{L}_t \cdot \frac{(1-\alpha_t)^{\sigma}\chi_t^{\sigma-1}}{
      \bigl[\alpha_t^{\sigma} R_{O,t}^{1-\sigma} +
      (1-\alpha_t)^{\sigma}(\chi_t R_{I,t})^{1-\sigma}\bigr]^{1/(1-\sigma)}\,
      R_{I,t}^{-\sigma}}.
      \label{eq:demand-I}
  \end{align}

  \textit{Step 3 (Spread determination: global existence and uniqueness).}
  The equilibrium spread $\mu_t^*$ is determined by the market-clearing
  condition~\eqref{eq:clearing} together with the sectoral demand
  system~\eqref{eq:net-demand}. Substituting~\eqref{eq:net-demand}
  into~\eqref{eq:clearing}, define the aggregate net-demand function
  \[
    F(\mu) \;\equiv\; \sum_s \mathfrak{q}_{s,m,t}(\mu)
    \;=\; B\mu + \sum_s \bigl[\alpha_{s,m} + \gamma_s X_{s,m,t}
    + \varepsilon_{s,m,t}\bigr],
  \]
  where $\mathcal{B} = \sum_s \beta_s < 0$ (by Remark~\ref{rem:B-negative}
  (maintained hypothesis), aggregate net demand is downward-sloping; this holds
  when demanders outweigh basis traders in aggregate price sensitivity).
  $F$ is a strictly decreasing affine function of $\mu$. \emph{Global
  existence:} as $\mu \to -\infty$, $F(\mu) \to +\infty$; as $\mu \to
  +\infty$, $F(\mu) \to -\infty$.
  By the intermediate value theorem applied to the continuous function $F$ on
  $\R$, there exists at least one $\mu_{m,t}^* \in \R$ such that $F(\mu_{m,t}^*)
  = 0$. \emph{Uniqueness:} since $\mathcal{B} < 0$, $F$ is strictly monotone
  decreasing, so the zero is unique. Solving explicitly:
  \[
    \mu_{m,t}^* = -\frac{1}{\mathcal{B}}\sum_s \bigl[\alpha_{s,m} +
    \gamma_s X_{s,m,t} + \varepsilon_{s,m,t}\bigr].
  \]

  \textit{Conclusion.} All six conditions in \cref{def:equilibrium} are
  satisfied. The tuple $\mathcal{E}_t$ is thus a well-defined equilibrium for
  any admissible parameter values.
  \hfill$\square$
\end{proof}

\begin{remark}[Uniqueness]
  The equilibrium spread $\mu_t^*$ is unique given $(\chi_t, B_t, U_t)$ because
  aggregate net demand is strictly monotone in price. The equilibrium quality
  $\chi_t^*$ is unique because it is pinned down by $p_t = \eta\theta_t^{-(1-\gamma)}$
  and the deterministic function $\chi$. Conditional on $B_t$ being exogenous,
  the full equilibrium tuple is unique. Multiplicity can arise when $B_t$ is
  endogenous (e.g., through $B_t^{HF}$ in \eqref{eq:HF-capacity}); see
  \cref{prop:run} and \cref{sec:extensions}.
\end{remark}

% ==============================================================================
\section{The Demand System}
\label{sec:demand-system}
% ==============================================================================

Following the logic of \citet{KoijenYogo2019}, we model each sector's signed
net USD position in a given market $m$ as a linear demand function. Let
$s \in \{FB, USbank, dealer, HF, AM, Corp\}$ index sectors and $m$ index
currency pair--tenor combinations.

\begin{definition}[Sectoral Net Demand]
  \label{def:net-demand}
  Sector $s$'s signed net USD position in market $m$ at date $t$ is
  \begin{equation}
    \mathfrak{q}_{s,m,t} = \alpha_{s,m} + \beta_s\,\mu_{m,t} +
    \gamma_s X_{s,m,t} + \varepsilon_{s,m,t},
    \label{eq:net-demand}
  \end{equation}
  where $\alpha_{s,m}$ is a sector--market fixed effect, $\beta_s$ is sector
  $s$'s price sensitivity (net demand slope with respect to the funding spread),
  $X_{s,m,t}$ is a vector of sector-specific demand shifters, and
  $\varepsilon_{s,m,t}$ is a mean-zero residual. For dollar demanders (FB,
  corporates, most AM): $\beta_s < 0$ (higher $\mu_{m,t}$ reduces demand). For
  basis traders (HF) acting as suppliers: $\beta_s > 0$ (higher $\mu_{m,t}$
  induces more supply).
\end{definition}

\begin{definition}[Equilibrium Price Equation]
  \label{def:price-eq}
  Imposing market clearing~\eqref{eq:clearing} and solving for $\mu_{m,t}$:
  \begin{equation}
    \mu_{m,t}^* = -\frac{1}{\mathcal{B}} \sum_s \bigl[\alpha_{s,m} +
    \gamma_s X_{s,m,t} + \varepsilon_{s,m,t}\bigr],
    \label{eq:price-eq}
  \end{equation}
  where $\mathcal{B} \equiv \sum_s \beta_s < 0$ is the aggregate net-demand
  slope (by Remark~\ref{rem:B-negative}). This requires $\mathcal{B} < 0$,
  i.e., that the sum of price sensitivities is negative: the aggregate net
  demand for synthetic dollars is downward-sloping in the funding spread. This
  holds when demanders outweigh basis traders in aggregate price-sensitivity.
\end{definition}

\begin{remark}[Price Impact Coefficient]
  \label{rem:price-impact}
  A sector-$s$ demand shock $\Delta X_{s,m,t}$ shifts $\mu_{m,t}$ by
  $-\gamma_s/\mathcal{B} > 0$ (since $\mathcal{B} < 0$ and $\gamma_{FB} > 0$
  for foreign bank demand shocks). This is the structural price impact
  coefficient. The IV estimator $\hat\beta = \hat\beta_{RF}/\hat\pi$ from the
  v8 design identifies this coefficient:
  $\hat\beta \xrightarrow{p} -\gamma_{FB}/\mathcal{B}$, assuming the
  shift-share instrument is valid (Assumption~\ref{ass:instrument} below).
\end{remark}

% ==============================================================================
\section{Main Results}
\label{sec:results}
% ==============================================================================

We now state and prove the seven main propositions. Each proposition is
self-contained: the assumptions it uses are stated explicitly at the head of the
proof. We work throughout with the static equilibrium of
\cref{def:equilibrium}.

% -----------------------------------------------------------------------
\subsection{Proposition 1: Demand and the Funding Spread}
% -----------------------------------------------------------------------

\begin{proposition}[Demand and Funding Spread]
  \label{prop:demand}
  Under Assumptions~\ref{ass:liq-req}--\ref{ass:spread} and
  \ref{ass:regularity}, at any interior equilibrium,
  \[
    \frac{\partial \mu_t^*}{\partial U_t} > 0.
  \]
  An increase in the mass of seekers raises the equilibrium funding spread,
  holding intermediary capacity $B_t$ fixed.
\end{proposition}

\begin{proof}
  The proof proceeds in two steps.

  \textit{Step 1 (Effect of $U_t$ on $\theta_t$ and $p_t$).}
  By \cref{def:matching}, $\theta_t = U_t/B_t$. Since $B_t$ is fixed,
  $\partial\theta_t/\partial U_t = 1/B_t > 0$. The matching probability is
  $p_t = \eta\,\theta_t^{-(1-\gamma)}$, so
  \[
    \frac{\partial p_t}{\partial U_t} =
    \frac{\partial p_t}{\partial \theta_t}\cdot\frac{\partial\theta_t}{\partial U_t}
    = -\eta(1-\gamma)\,\theta_t^{-(2-\gamma)} \cdot \frac{1}{B_t} < 0.
  \]
  By Assumption~\ref{ass:quality}, $\chi_t = \chi(p_t)$ is increasing in $p_t$,
  so $\partial\chi_t/\partial U_t = \chi'(p_t)\cdot\partial p_t/\partial U_t < 0$.

  \textit{Step 2 (Effect of $\chi_t$ on $\mu_t^*$: cross-layer transmission).}
  The key challenge is that $\chi_t$ enters the CES aggregator~\eqref{eq:CES}
  but, as written, does not appear explicitly as an argument of the sectoral
  demand functions $\mathfrak{q}_{s,m,t}(\mu_{m,t})$ in~\eqref{eq:net-demand}.
  We now show how the quality parameter transmits from the CES layer to the
  demand-system layer through the effective cost of inside liquidity.

  \textbf{Quality-adjusted price.} In the CES aggregator, one unit of inside
  liquidity $I_t$ delivers $\chi_t$ units of effective liquidity service.
  Accordingly, the cost-minimization program~\eqref{eq:cost-min} is equivalent
  to minimizing $R_{O,t} O_t + R_{I,t} I_t$ subject to achieving $\bar{L}_t$
  units of CES service. The cost of obtaining one unit of \emph{effective}
  service via the inside channel is $R_{I,t}/\chi_t$, not $R_{I,t}$. Define
  the \emph{quality-adjusted price of inside liquidity}:
  \[
    \tilde{R}_{I,t} \;\equiv\; \frac{R_{I,t}}{\chi_t},
  \]
  and the corresponding \emph{quality-adjusted spread}:
  \[
    \tilde{\mu}_t \;\equiv\; \tilde{R}_{I,t} - R_{O,t}
    \;=\; \frac{R_{I,t}}{\chi_t} - R_{O,t}.
  \]
  In terms of $\tilde{\mu}_t$, the demand-system specification~\eqref{eq:net-demand}
  correctly represents sectoral net demand: sector $s$ responds to the
  quality-adjusted spread because it compares the effective cost-per-service-unit
  of inside liquidity to that of outside liquidity. We may therefore write
  $\mathfrak{q}_{s,m,t} = \alpha_{s,m} + \beta_s \tilde{\mu}_{m,t} +
  \gamma_s X_{s,m,t} + \varepsilon_{s,m,t}$.

  \textbf{Equilibrium quality-adjusted spread.} Imposing market
  clearing~\eqref{eq:clearing}:
  \[
    \tilde{\mu}_{m,t}^* = -\frac{1}{\mathcal{B}}\sum_s \bigl[\alpha_{s,m} +
    \gamma_s X_{s,m,t} + \varepsilon_{s,m,t}\bigr],
  \]
  where $\mathcal{B} = \sum_s \beta_s < 0$ (by Remark~\ref{rem:B-negative}).
  Note that $\tilde{\mu}^*$ does not depend on $\chi_t$: the equilibrium
  quality-adjusted spread is determined solely by the supply-demand balance in
  terms of effective service units.

  \textbf{Deriving $\partial\mu_t^*/\partial\chi_t$ via the implicit function theorem.}
  \emph{Note: the price-recovery identity $\mu_t = \chi_t(\tilde\mu_t^* + R_{O,t})
  - R_{O,t}$ does not yield the equilibrium comparative static
  $\partial\mu_t^*/\partial\chi_t$ because $\tilde\mu_t^*$ is defined conditional
  on a fixed $\chi_t$; the correct approach uses the implicit function theorem on
  the market-clearing condition.}

  The correct argument proceeds as follows. Consider the market-clearing condition
  written as $F(\mu;\chi_t) \equiv \sum_s \mathfrak{q}_{s,m,t}(\mu,\chi_t) = 0$.
  Holding the posted rate $\mu$ fixed, a fall in $\chi_t$ reduces the effective
  service of $I_t$, so agents must demand more $I_t$ to satisfy $\bar{L}_t$:
  this raises $\mathfrak{q}_{FB,m,t}$ (more dollar-taking) at any given $\mu$.
  Hence $\partial\mathfrak{q}_{FB,m,t}/\partial\chi_t < 0$ (quality fall raises
  quantity demanded of inside liquidity). Summing across sectors:
  $\partial F(\mu;\chi_t)/\partial\chi_t < 0$ (the aggregate demand curve shifts
  up --- more inside liquidity demanded at any given $\mu$, requiring a higher
  $\mu$ to restore market clearing). By the implicit function theorem:
  \[
    \frac{\partial\mu_t^*}{\partial\chi_t}
    = -\frac{\partial F/\partial\chi_t}{\partial F/\partial\mu}
    = -\frac{(\text{negative})}{\mathcal{B}} = -\frac{(\text{negative})}{(\text{negative})}
    < 0.
  \]
  Hence $\partial\mu_t^*/\partial\chi_t < 0$: quality deterioration raises
  the equilibrium posted spread. Combining:
  \[
    \frac{\partial\mu_t^*}{\partial U_t}
    = \underbrace{\frac{\partial\mu_t^*}{\partial\chi_t}}_{<\,0}
      \cdot\underbrace{\chi'(p_t)}_{>\,0}
      \cdot\underbrace{\frac{\partial p_t}{\partial U_t}}_{<\,0}
    \;>\; 0. \qquad \checkmark
  \]

  \textbf{Note on propagation to Props 2, 4, and 5.} The same
  $\chi_t \to \mu_t^*$ transmission — quality deterioration raises the
  equilibrium posted spread through the demand-quantity shift and the IFT on
  market clearing — underlies \cref{prop:capacity}, \cref{prop:qe}, and
  \cref{prop:swaplines}. In each case, the capacity change affects $\chi_t$
  via $p_t$ via $\theta_t$, and the same IFT argument applies with the sign
  reversed where $\chi_t$ rises (capacity increases) rather than falls.

  \textit{Assumptions used:} Assumption~\ref{ass:liq-req} ensures $\bar{L}_t >
  0$ (binding constraint); Assumption~\ref{ass:sigma} ensures the interior
  solution; Assumption~\ref{ass:quality} delivers $\chi'(p) > 0$;
  Assumption~\ref{ass:regularity} ensures differentiability.
  \hfill$\square$
\end{proof}

\begin{corollary}[IV Estimand]
  \label{cor:IV}
  The causal parameter $\partial\mu_t^*/\partial U_t$ is identified by a valid
  instrument $Z_t$ that shifts $U_t$ (or its empirical proxy $Q_{FB,m,t}$) but
  is excluded from the spread equation. The shift-share instrument
  $Z^{MMF}_{m,t} = \mathrm{Exposure}^{MMF}_{m,t-1} \times \mathrm{FundingShock}_t$
  satisfies these conditions under the additional assumption that pre-determined
  exposure weights $w_{i,m}^{pre}$ are uncorrelated with unobserved determinants
  of $\mu_{m,t}$ other than through $Q_{FB,m,t}$ (see
  Assumption~\ref{ass:instrument} below).
\end{corollary}

\begin{assumption}[Shift-Share Instrument Validity]
  \label{ass:instrument}
  The shift-share instrument $Z^{MMF}_{m,t}$ satisfies: (i) relevance:
  $\mathrm{Cov}(Z^{MMF}_{m,t}, Q_{FB,m,t}) \neq 0$; (ii) exogeneity of the
  common shock: $\mathrm{FundingShock}_t$ is orthogonal to market-specific
  $\varepsilon_{m,t}$ conditional on time fixed effects; (iii) exogeneity of
  the shares: the pre-determined exposure weights $w_{i,m}^{pre}$ are
  orthogonal to $\varepsilon_{m,t}$ conditional on market fixed effects.
  Condition (iii) requires that dominant-share banks do not systematically
  differ from low-exposure banks in their unobserved FX swap price-setting
  ability --- a testable condition via balance tests on exposure-weighted bank
  characteristics.
\end{assumption}

\begin{remark}[Goldsmith-Pinkham et al.\ Critique]
  Assumption~\ref{ass:instrument}(iii) is stronger than mere pre-determination
  of the weights. Following \citet{GoldsmithPinkham2020}, validity requires
  that the dominant-share banks are not systematically different from others in
  ways that independently predict $\mu_{m,t}$. This must be verified empirically
  via balance tests before the IV estimates can be given a causal interpretation.
\end{remark}

% -----------------------------------------------------------------------
\subsection{Proposition 2: Capacity and the Funding Spread}
% -----------------------------------------------------------------------

\begin{proposition}[Capacity and Funding Spread]
  \label{prop:capacity}
  Under Assumptions~\ref{ass:liq-req}--\ref{ass:regularity}, at any interior
  equilibrium,
  \[
    \frac{\partial \mu_t^*}{\partial B_t} < 0.
  \]
  An increase in intermediary capacity reduces the equilibrium funding spread,
  holding seeker mass $U_t$ fixed.
\end{proposition}

\begin{proof}
  By \cref{def:matching}, $\partial\theta_t/\partial B_t = -U_t/B_t^2 < 0$, so
  market tightness falls when capacity rises. By the same chain as in the proof
  of \cref{prop:demand}:
  \[
    \frac{\partial p_t}{\partial B_t} =
    \frac{\partial p_t}{\partial\theta_t}\cdot\frac{\partial\theta_t}{\partial B_t}
    = \underbrace{-\eta(1-\gamma)\theta_t^{-(2-\gamma)}}_{<0}
      \cdot\underbrace{\left(-\frac{U_t}{B_t^2}\right)}_{<0} > 0.
  \]
  Hence $\partial\chi_t/\partial B_t = \chi'(p_t)\cdot\partial p_t/\partial
  B_t > 0$: higher capacity raises quality.

  Following the same logic as \cref{prop:demand} Step 2 with the sign reversed:
  higher $\chi_t$ raises effective inside liquidity supply, relaxing the
  scarcity of conversion capacity and reducing the spread. Formally,
  $\partial\mu_t^*/\partial B_t = (\partial\mu_t^*/\partial\chi_t)\cdot
  \chi'(p_t)\cdot(\partial p_t/\partial B_t)$. Since $\partial\mu_t^*/
  \partial\chi_t < 0$ and $\chi'(p_t)\cdot\partial p_t/\partial B_t > 0$,
  the product is negative.
  \hfill$\square$
\end{proof}

\begin{remark}[Causal Identification of \cref{prop:capacity}]
  Unlike \cref{prop:demand}, the capacity channel cannot be identified by a
  clean instrument in the current design, because $B_t$ is endogenous ---
  dealer constraints tighten precisely in high-stress states. The empirical
  strategy (v8 P6--P7) tests the prediction indirectly via interaction terms
  $\hat{Q}_{FB,m,t} \times \mathrm{DealerConstraint}_t$ and $\hat{Q}_{FB,m,t}
  \times Q_t^{QE}$, where $\mathrm{DealerConstraint}_t$ is a regime indicator.
  This identifies heterogeneity in price impact across capacity regimes, not
  the causal effect of $B_t$ directly. The mapping from \cref{prop:capacity} to
  the empirical test is therefore $\partial\mu_t^*/\partial B_t < 0$ manifesting
  as a larger price-impact coefficient $\beta$ when $B_t$ is constrained:
  $\beta|_{B_t\text{ low}} > \beta|_{B_t\text{ high}}$.
\end{remark}

% -----------------------------------------------------------------------
\subsection{Proposition 3: Non-linearity and Threshold Effects}
% -----------------------------------------------------------------------

\begin{proposition}[Non-linearity: Threshold in Market Tightness]
  \label{prop:nonlinear}
  Under Assumptions~\ref{ass:liq-req}--\ref{ass:regularity}, and additionally
  assuming (i) $\chi''(p) > 0$ (quality is strictly convex in matching
  probability) and (ii) $\Phi''(\chi) > 0$ (the equilibrium spread is strictly
  convex in quality, where $\Phi\colon(0,1]\to\R$ is the mapping from quality
  to spread derived from the equilibrium price equation), the second derivative
  of the equilibrium spread with respect to market tightness satisfies:
  \[
    \frac{\partial^2 \mu_t^*}{\partial \theta_t^2} > 0,
  \]
  so the funding spread is a convex function of market tightness. Under
  additional regularity, there exists a threshold $\theta^* > 0$ such that
  for $\theta_t > \theta^*$ the rate of increase $\partial\mu_t^*/\partial
  \theta_t$ exceeds the rate for $\theta_t < \theta^*$, delivering a locally
  sharp escalation in the spread.
\end{proposition>

\begin{proof}
  Write $\mu_t^* = \Phi(\chi(p(\theta_t)))$ where $\Phi$ captures the mapping
  from quality to spread (from the equilibrium price equation), $p(\theta) =
  \eta\,\theta^{-(1-\gamma)}$, and $\Phi$ is decreasing in $\chi$ (higher
  quality reduces the spread). Then:
  \[
    \frac{d\mu_t^*}{d\theta_t} = \Phi'(\chi)\cdot\chi'(p)\cdot p'(\theta) > 0,
  \]
  using $\Phi' < 0$, $\chi' > 0$, and $p'(\theta) < 0$. The second derivative
  is:
  \begin{align*}
    \frac{d^2\mu_t^*}{d\theta_t^2} &=
    \Phi''(\chi)\bigl[\chi'(p)p'(\theta)\bigr]^2
    + \Phi'(\chi)\chi''(p)\bigl[p'(\theta)\bigr]^2
    + \Phi'(\chi)\chi'(p)p''(\theta).
  \end{align*}
  The term $p''(\theta) = \eta(1-\gamma)(2-\gamma)\theta^{-(3-\gamma)} > 0$
  (concave $p$ in $B$ but convex in $\theta$). Expanding each of the three
  terms in $d^2\mu^*/d\theta^2$:
  \begin{itemize}[noitemsep]
    \item \textbf{Term 1}: $\Phi''(\chi)[\chi'(p)p'(\theta)]^2 > 0$ under
    $\Phi'' > 0$ (the maintained condition (ii)).
    \item \textbf{Term 2}: $\Phi'(\chi)\chi''(p)[p'(\theta)]^2$. Since
    $\Phi' < 0$ and $\chi'' > 0$ (condition (i)) and $[p']^2 > 0$, this term
    is negative.
    \item \textbf{Term 3}: $\Phi'(\chi)\chi'(p)p''(\theta)$. Since $\Phi'<0$,
    $\chi'>0$, and $p''>0$, this term is also negative.
  \end{itemize}
  The first term dominates when $\Phi''$ is sufficiently large relative to
  $|\Phi'|$: formally, a sufficient condition is
  \[
    \Phi''(\chi) > |\Phi'(\chi)| \cdot
    \frac{\chi''(p)[p'(\theta)]^2 + \chi'(p)p''(\theta)}{[\chi'(p)p'(\theta)]^2}.
  \]
  Both conditions (i) and (ii) together are therefore jointly sufficient for
  $d^2\mu^*/d\theta^2 > 0$. Neither alone is sufficient in general: if
  $\Phi'' = 0$ (linear $\Phi$), conditions (i) and matching-function convexity
  alone determine the sign, which can go either way depending on the relative
  magnitudes of Terms 2 and 3.

  The threshold $\theta^*$ arises naturally under a globally convex $\mu^*(\theta)$
  as any value above which the rate of increase exceeds the baseline; or, if
  $\chi$ is S-shaped (concave then convex), as the inflection point of
  $\mu^*(\theta)$, which is computable in closed form under $\chi(p) = p^\delta$
  with $\delta > 1$ (in which case global convexity holds).

  \textit{Assumptions used:} Assumptions~\ref{ass:quality} (differentiability),
  \ref{ass:matching} (matching function), \ref{ass:regularity} (differentiability
  of equilibrium). The non-linearity result additionally requires both $\chi''
  > 0$ and $\Phi'' > 0$ as stated in the proposition; neither is inherited from
  the standing assumptions.
  \hfill$\square$
\end{proof}

\begin{remark}[Sufficient Conditions for $\Phi'' > 0$ and $\chi'' > 0$]
  \label{rem:convexity-conditions}
  The condition $\Phi'' > 0$ holds when the equilibrium spread responds more
  than proportionally to quality deterioration. This is implied by the
  demand-system formulation when $\mathcal{B} = \sum_s\beta_s$ is small in
  absolute value (thin market): the price impact of any demand shift is
  $-\gamma_s/\mathcal{B}$, which grows as $|\mathcal{B}| \to 0$; if quality
  deterioration is concurrently associated with a reduction in $|\mathcal{B}|$
  (fewer active participants in stress), then $\Phi$ is convex in $\chi$. Under CES aggregation with
  $\sigma > 1$, a deterioration in $\chi$ also amplifies the effective price
  gap between inside and outside liquidity in a superlinear manner, providing
  a structural micro-foundation for $\Phi'' > 0$.

  The condition $\chi'' > 0$ requires that quality deteriorates convexly as
  matching becomes harder. This holds under the parametric family
  $\chi(p) = p^\delta$ with $\delta > 1$. However, $\chi'' > 0$ is not
  empirically unambiguous: if $\chi(p) = \chi_{\min} + (1-\chi_{\min})p$ is
  linear, or if $\chi$ is concave ($\delta < 1$), then $\chi'' \leq 0$ and
  this result fails. Failure of the joint condition produces $d^2\mu^*/d\theta^2
  \leq 0$ (concavity or linearity), eliminating the threshold effect. The paper
  should present \cref{prop:nonlinear} as conditional on the convex subfamily
  of quality functions and note this restriction explicitly.
\end{remark}

\begin{remark}[Empirical Testability of \cref{prop:nonlinear}]
  The threshold effect in \cref{prop:nonlinear} requires a continuous
  $\hat\theta_{m,t} = Q_{FB,m,t}/\mathrm{DealerCapacity}_t$ series and a
  structural break or threshold regression. The explorer feasibility report
  (68/100) grades $\theta_t$ measurement at C and classifies this proposition
  as ``not testable with the current design.'' The interaction specifications
  in v8 (interacting $Q_{FB}$ with a DealerConstraint binary) test
  monotone heterogeneity, not a threshold. A formal threshold test requires
  a continuous $B_t$ series and an econometric threshold-regression procedure
  (e.g., \citealt{Hansen2000}). This should be stated clearly in the paper as
  an open empirical question.
\end{remark}

% -----------------------------------------------------------------------
\subsection{Proposition 4: Quarter-End Regulatory Amplification}
% -----------------------------------------------------------------------

\begin{proposition}[Quarter-End Amplification]
  \label{prop:qe}
  Under Assumptions~\ref{ass:liq-req}--\ref{ass:supply} and
  \ref{ass:regularity}, if $\omega > 0$ in \eqref{eq:FBO-capacity}, then at
  quarter-end dates ($Q_t = 1$):
  \[
    B_t|_{Q_t=1} < B_t|_{Q_t=0} \implies \mu_t^*|_{Q_t=1} >
    \mu_t^*|_{Q_t=0},
  \]
  holding all other exogenous variables constant. Furthermore, the price impact
  of a given demand shock $\Delta U_t$ is larger at quarter-end:
  \[
    \left.\frac{\partial \mu_t^*}{\partial U_t}\right|_{B_t^{QE}}
    > \left.\frac{\partial \mu_t^*}{\partial U_t}\right|_{B_t^{normal}},
  \]
  where $B_t^{QE} < B_t^{normal}$.
\end{proposition}

\begin{proof}
  The first claim is a direct application of \cref{prop:capacity} with
  $\Delta B_t = -\omega < 0$: a reduction in FBO capacity at quarter-end
  raises $\theta_t$ (since $\partial\theta_t/\partial B_t < 0$), lowers $p_t$
  and $\chi_t$, and raises $\mu_t^*$.

  For the price-impact amplification: the partial derivative
  $\partial\mu_t^*/\partial U_t = \Phi'(\chi)\chi'(p)\cdot\partial p/\partial U_t$.
  Since $\partial p/\partial U_t = -\eta(1-\gamma)\theta_t^{-(2-\gamma)}/B_t$,
  the magnitude is increasing in $\theta_t^{-(2-\gamma)}/B_t$, which is
  larger when $B_t$ is smaller (quarter-end). This establishes the
  interaction prediction.

  \textit{Attribution caveat:} This proof attributes the quarter-end effect
  entirely to the supply channel ($B_t \downarrow$). The empirical test
  (v8 P7) identifies the reduced-form effect ($\mu_t \uparrow$ at quarter-end)
  but cannot separately identify the supply-side channel from a simultaneous
  demand-side effect ($U_t \uparrow$ from portfolio rebalancing at quarter-end).
  Both channels predict $\mu_t \uparrow$ and are consistent with the proposition.
  A sign test between the two channels would require a separate instrument for
  supply-side variation at quarter-end.
  \hfill$\square$
\end{proof}

% -----------------------------------------------------------------------
\subsection{Proposition 5: Central Bank Swap Lines}
% -----------------------------------------------------------------------

\begin{proposition}[Swap Lines and Spread Compression]
  \label{prop:swaplines}
  Under Assumptions~\ref{ass:liq-req}--\ref{ass:supply} and
  \ref{ass:regularity}, an increase in central-bank swap-line capacity
  $\Delta B_t^{CBswap} > 0$ raises $B_t$ and reduces $\mu_t^*$:
  \[
    \frac{\partial \mu_t^*}{\partial B_t^{CBswap}} < 0.
  \]
  The reduction is strictly larger when market tightness $\theta_t$ is high
  (i.e., in stressed states), because $|\partial\mu_t^*/\partial B_t|$ is
  increasing in $\theta_t$ when $\chi$ is convex (by \cref{prop:nonlinear}).
\end{proposition}

\begin{proof}
  By \cref{def:capacity}, $B_t^{CBswap}$ enters $B_t$ additively and with
  coefficient 1. Therefore $\partial B_t/\partial B_t^{CBswap} = 1$, and the
  result follows directly from \cref{prop:capacity}:
  $\partial\mu_t^*/\partial B_t^{CBswap} = \partial\mu_t^*/\partial B_t < 0$.

  The stress-state amplification follows from \cref{prop:nonlinear}: when
  $\theta_t$ is high, each unit increase in $B_t$ produces a larger reduction
  in $\theta_t$ (in proportional terms, since $\partial\theta_t/\partial B_t
  = -U_t/B_t^2$ is larger in magnitude when $B_t$ is small), and therefore
  a larger improvement in $p_t$ and $\chi_t$, and a larger reduction in
  $\mu_t^*$.

  \textit{What assumptions are needed.} The chain $B_t^{CBswap} \uparrow \to
  p_t \uparrow \to \chi_t \uparrow \to \mu_t^* \downarrow$ requires:
  (i) Assumption~\ref{ass:quality} for $\chi'(p) > 0$ (quality rises with
  matching probability); (ii) \cref{prop:capacity} for $\partial\mu_t^*/
  \partial\chi_t < 0$; (iii) that swap-line dollars are accepted by the market
  as equivalent to private outside liquidity (i.e., $B_t^{CBswap}$ enters
  additively in $B_t$, not at a discount). If this last condition fails ---
  e.g., because swap-line dollars carry stigma or are available only to
  specific counterparties --- the effective capacity increase is $\delta_t
  B_t^{CBswap}$ with $\delta_t < 1$, and the chain still holds but with a
  smaller magnitude.
  \hfill$\square$
\end{proof}

\begin{remark}[Identification Challenge for \cref{prop:swaplines}]
  Swap lines are activated precisely in crises, so any event-study estimate of
  $\partial\mu_t^*/\partial B_t^{CBswap}$ conflates the policy effect with
  natural crisis resolution. Cross-currency variation (currency pairs covered
  vs.\ not covered by swap lines) provides the cleanest identification, but
  requires assuming that uncovered pairs are valid counterfactuals. This is
  plausible for close substitutes (e.g., EUR/USD vs.\ NOK/USD during
  March 2020) but less so for fundamentally different pairs.
\end{remark}

% -----------------------------------------------------------------------
\subsection{Proposition 6: Intermediary Heterogeneity}
% -----------------------------------------------------------------------

\begin{proposition}[Provider Heterogeneity]
  \label{prop:heterogeneity}
  Under Assumption~\ref{ass:supply}, the three main intermediary types exhibit
  qualitatively different supply behaviors:
  \begin{enumerate}[label=(\alph*)]
    \item \textbf{FBO branches and primary dealers} supply capacity $B_t^{FBO}$
    and $B_t^{dealer}$ that is (weakly) decreasing in stress $S_t$ and in the
    quarter-end indicator, i.e., pro-cyclically constrained.
    \item \textbf{Hedge funds} supply capacity $B_t^{HF} = \phi\max(\mu_t -
    c^{HF},0)\kappa_t$ that is increasing in $\mu_t$ when $\mu_t > c^{HF}$
    (they enter when the basis is profitable) and zero when $\mu_t \leq c^{HF}$
    or $\kappa_t = 0$ (they exit in extreme stress). In particular, $B_t^{HF}$
    is non-monotone in stress: it is increasing when $\mu_t$ rises mildly
    (basis-trade entry) and falling when $\kappa_t$ collapses in acute stress.
    \item \textbf{Central-bank swap lines} supply $B_t^{CBswap}$ that is
    explicitly counter-cyclical (deployed in crises by policy design).
  \end{enumerate}
\end{proposition}

\begin{remark}[Maintained Hypothesis on Price-Sensitivity Ordering]
  \label{rem:beta-ordering}
  The proposition as stated does not include a claim about $|\beta_{FB}|$
  versus $\beta_{HF}$. This ordering is \emph{not} derived from model
  primitives: the coefficients $\beta_s$ are free parameters of the linear
  demand system~\eqref{eq:net-demand} that are estimated from data, not
  pinned down by the theoretical structure. We therefore state the ordering
  as a maintained empirical hypothesis rather than a theoretical result.

  \textbf{Maintained Hypothesis (Price-Sensitivity Ordering):}
  We impose the empirically motivated ordering
  \[
    |\beta_{FB}| < \beta_{HF},
  \]
  reflecting the structural interpretation that foreign banks face inelastic
  payment obligations ($\bar{L}_t$ predetermined by Assumption~\ref{ass:liq-req})
  while hedge funds face no analogous inelastic demand and enter only when the
  basis exceeds their entry cost $c^{HF}$. This ordering is consistent with
  the model's demand structure but is not derived from it. It is tested
  empirically in the demand system estimation (v8 Part 5, Prediction P3): the
  null hypothesis $|\beta_{FB}| \geq \beta_{HF}$ is rejected if and only if
  foreign-bank demand is indeed more price-inelastic than HF supply.
  Failure to reject would indicate that the structural interpretation of FBO
  inelasticity is not supported by the data.
\end{remark}

\begin{proof}[Proof of \cref{prop:heterogeneity}]
  Parts (a) and (c) follow directly from Assumption~\ref{ass:supply}(a) and
  (c): FBO capacity decreases in $S_t$ (stress) and $Q_t = 1$ (quarter-end)
  by the parametric structure of \eqref{eq:FBO-capacity}; central-bank swap
  lines are counter-cyclical by policy design.

  For part (b): the hedge fund supply function $B_t^{HF} = \phi
  \max(\mu_t - c^{HF}, 0)\kappa_t$ has $\partial B_t^{HF}/\partial\mu_t =
  \phi\kappa_t > 0$ when $\mu_t > c^{HF}$, and $= 0$ otherwise. When
  $\kappa_t \to 0$ (acute stress), $B_t^{HF} \to 0$ regardless of $\mu_t$.
  This creates the non-monotone pattern: moderate stress (rising $\mu_t$,
  intact $\kappa_t$) increases HF supply, while severe stress ($\kappa_t \to 0$)
  withdraws it entirely.

  The price-sensitivity ordering $|\beta_{FB}| < \beta_{HF}$ is not a derived
  result; it is a maintained empirical hypothesis stated in \cref{rem:beta-ordering}.
  \hfill$\square$
\end{proof}

% -----------------------------------------------------------------------
\subsection{Proposition 7: Run Dynamics and Amplification (New)}
% -----------------------------------------------------------------------

The preceding propositions characterize the \emph{comparative statics} of a
unique interior equilibrium. We now examine whether the amplification mechanism
in Section~11.1 of the original draft can produce \emph{multiple equilibria},
one of which corresponds to a ``run'' state.

\begin{definition}[Reduced-Form Equilibrium Map]
  \label{def:eqmap}
  Define the \emph{equilibrium correspondence} $\Psi\colon \R_{+} \to \R_{+}$
  by
  \begin{equation}
    \mu^* = \Psi(\mu^*;\, U_0, B_0, \alpha, \bar{L}, \kappa),
    \label{eq:fixed-point}
  \end{equation}
  where $\Psi(\mu; \cdot)$ is the value of the equilibrium spread that solves
  \cref{def:equilibrium} when $B_t^{HF}$ is evaluated at $\mu$ (using
  \eqref{eq:HF-capacity}) and all other components of $B_t$ are fixed. The
  fixed points of $\Psi$ are the equilibria of the full model with endogenous
  HF supply.
\end{definition}

\begin{proposition}[Run Equilibrium: Conditions for Multiple Equilibria]
  \label{prop:run}
  Suppose hedge fund capacity is endogenous as in Assumption~\ref{ass:supply}(b),
  with $\phi > 0$ and $\kappa_t = \kappa > 0$ constant. Define
  \[
    \Lambda(\mu) \equiv \frac{\partial\Psi(\mu;\cdot)}{\partial\mu}
    = \frac{\partial\mu^*}{\partial B_t^{HF}}
      \cdot \phi\kappa\,\mathds{1}[\mu > c^{HF}].
  \]
  \begin{enumerate}[label=(\alph*)]
    \item \textbf{(Unique equilibrium, normal times):} If $|\Lambda(\mu)| < 1$
    for all $\mu \in [c^{HF}, \bar\mu]$, where $\bar\mu$ is the spread in
    the absence of HF participation, the fixed-point equation~\eqref{eq:fixed-point}
    has a unique solution $\mu^{**}$ strictly below $\bar\mu$. HF entry
    dampens the spread.

    \item \textbf{(Multiple equilibria, run state):} Suppose
    $|\Lambda(\mu)| > 1$ for all $\mu \in (c^{HF}, \bar\mu)$. Then the
    fixed-point map $\Psi$ is not a contraction on $(c^{HF}, \bar\mu)$,
    and the model admits (at least) two stable equilibria: a \emph{normal
    equilibrium} $\mu^{low}$ with active HF participation and a \emph{run
    equilibrium} $\mu^{high} > \mu^{low}$ with minimal HF participation.

    \item \textbf{(Amplification through quality):} When $\kappa$ falls
    exogenously (acute stress withdrawal), $B_t^{HF}$ falls at any fixed
    $\mu$, which raises $\theta_t$, lowers $p_t$ and $\chi_t$, and raises
    $\mu^*$. This triggers further HF withdrawal at the new higher $\mu$ if
    the entry condition $\mu > c^{HF}$ is undermined by risk-capacity
    exhaustion. The amplification loop is:
    \[
      \kappa \downarrow \;\to\; B_t^{HF} \downarrow \;\to\; \theta_t \uparrow
      \;\to\; p_t \downarrow \;\to\; \chi_t \downarrow \;\to\; \mu_t^* \uparrow
      \;\to\; \text{demand for conversion} \uparrow \;\to\; U_t \uparrow
      \;\to\; \theta_t \uparrow\uparrow.
    \]
    This self-reinforcing loop constitutes a \emph{run on inside synthetic dollar
    liquidity}: the deterioration of conversion quality drives agents to demand
    more conversion, which further degrades quality.
  \end{enumerate}
\end{proposition}

\begin{proof}
  \textit{Part (a).} Define $F(\mu) \equiv \Psi(\mu;\cdot) - \mu$. By
  \cref{prop:capacity}, $\partial\mu^*/\partial B_t^{HF} < 0$ (higher capacity
  lowers the spread). The HF supply function has $\partial B_t^{HF}/\partial\mu
  = \phi\kappa > 0$. Therefore $\partial\Psi/\partial\mu = \Lambda(\mu) < 0$
  for $\mu > c^{HF}$: HF entry is a stabilizing force (higher $\mu$ induces
  more supply, which reduces $\mu^*$). When $|\Lambda(\mu)| < 1$, the Banach
  fixed-point theorem applied to $F$ on $[c^{HF}, \bar\mu]$ guarantees a unique
  fixed point, since $F$ is a contraction. The fixed point $\mu^{**}$ is the
  unique equilibrium.

  \textit{Part (b).} We establish multiple equilibria using the intermediate
  value theorem applied to $F(\mu) = \Psi(\mu) - \mu$. We first establish
  the correct boundary conditions for $F$.

  \textbf{Boundary condition at $\mu \to c^{HF}+$.} As $\mu \downarrow c^{HF}$,
  $B_t^{HF} = \phi\max(\mu - c^{HF},0)\kappa \to 0$, so HF supply vanishes.
  The market then operates as if hedge funds are absent, and the equilibrium
  spread approaches the no-HF value $\bar\mu$ (the fixed point under
  $B_t^{HF} = 0$). By assumption $\bar\mu > c^{HF}$ (otherwise HF never
  enter). Thus $\Psi(c^{HF}+) = \bar\mu > c^{HF}$, giving
  \[
    F(c^{HF}+) = \Psi(c^{HF}+) - c^{HF} = \bar\mu - c^{HF} > 0.
  \]

  \textbf{Boundary condition as $\mu \to \infty$.} As $\mu \to \infty$,
  $B_t^{HF} = \phi(\mu - c^{HF})\kappa \to \infty$. By \cref{prop:capacity},
  $\mu^* \to 0$ as HF capacity grows without bound (infinite capacity eliminates
  the spread). Thus $\Psi(\mu) \to 0$ as $\mu \to \infty$, giving
  \[
    F(\mu) = \Psi(\mu) - \mu \to 0 - \infty = -\infty < 0.
  \]

  \textbf{Existence of at least one fixed point.} Since $F(c^{HF}+) > 0$ and
  $F(\mu) < 0$ for all sufficiently large $\mu$, by the intermediate value
  theorem (applied to the continuous function $F$ on $(c^{HF}, \infty)$) there
  exists at least one $\mu_1^* \in (c^{HF}, \infty)$ with $F(\mu_1^*) = 0$.

  \textbf{Existence of a second fixed point under the global condition
  $|\Lambda(\mu)| > 1$ on $(c^{HF}, \mu_1^*)$.}
  Assume $|\Lambda(\mu)| > 1$ (i.e., $F'(\mu) = \Lambda(\mu) - 1 < -1$)
  for all $\mu \in (c^{HF}, \mu_1^*)$ --- a sufficient condition for the
  existence of multiple fixed points. Under this global condition, $F$ has
  slope strictly less than $-1$ throughout the interval. Starting from
  $F(c^{HF}+) = \bar\mu - c^{HF} > 0$, after a displacement of
  $\Delta \equiv \bar\mu - c^{HF}$ from the left boundary, $F$ descends
  by more than $\Delta$ (since slope $< -1$ throughout). Therefore $F$
  crosses zero before reaching $\mu_1^*$, establishing $F(\mu^{low}) = 0$
  for some $\mu^{low} \in (c^{HF}, \mu_1^*)$. The two fixed points are
  $\mu^{low}$ and $\mu_1^*$. The lower fixed point $\mu^{low}$ corresponds
  to the normal equilibrium (HF actively stabilizing); the upper fixed point
  $\mu^{high} \equiv \mu_1^*$ corresponds to the run equilibrium (spread
  at the no-stabilization level). The sign of $F'$ at each zero can be
  verified from the global slope condition: $F'(\mu^{low}) > 0$ (stable)
  and $F'(\mu^{high}) < -1 < 0$ (stable), consistent with two stable
  equilibria separated by a region where $F < 0$.

  \textit{Sufficient condition for $|\Lambda| > 1$.} Since $\Lambda(\mu) =
  (\partial\mu^*/\partial B_t)\cdot\phi\kappa$, and $|\partial\mu^*/\partial B_t|
  = (1/|\mathcal{B}|)\cdot(\partial\mu^*/\partial\chi)\cdot\chi'(p)\cdot|\partial p/
  \partial B_t|$, the condition $|\Lambda| > 1$ is more likely when: (i) the
  aggregate price sensitivity $|\mathcal{B}|$ is small (thin market); (ii) $\phi\kappa$
  is large (high HF capacity and sensitivity); or (iii) $\chi'(p)$ is large
  (quality is highly sensitive to matching probability, i.e., the conversion
  market is fragile). Combining (i) and (iii) corresponds to a fragile market
  with few alternative suppliers --- precisely the scenario in acute dollar
  funding stress.

  \begin{remark}[Global vs.\ Local Amplification Condition]
    \label{rem:global-lambda}
    The condition $|\Lambda(\mu)| > 1$ for all $\mu \in (c^{HF}, \mu_1^*)$ is
    sufficient but not necessary for the existence of two fixed points. A weaker
    condition --- local $|\Lambda(\hat\mu)| > 1$ at some interior point $\hat\mu$
    combined with a bound on $F(c^{HF}+) = \bar\mu - c^{HF}$ --- would suffice
    but requires case analysis of the trajectory of $F$ between $c^{HF}$ and
    $\hat\mu$. The global condition avoids this case analysis and is empirically
    meaningful: in stressed markets, the amplification loop is strong throughout
    the entire tightness range, not merely at an isolated point.
  \end{remark}

  \textit{Part (c).} The amplification loop follows from the chain rule applied
  to $\Psi$. An exogenous fall in $\kappa$ reduces $B_t^{HF}$ at fixed $\mu$,
  raising $\theta_t$ and lowering $p_t$ (by \cref{def:matching}). Lower $p_t$
  lowers $\chi_t$ (by \cref{def:quality}). Lower $\chi_t$ raises $\mu_t^*$
  (by the equilibrium price equation, since effective inside liquidity supply
  has fallen). Higher $\mu_t^*$ increases the demand for conversion: agents
  facing higher conversion costs demand more conversions (they must pay more to
  access safe dollars, so they need more of them to meet $\bar{L}_t$). This
  raises $U_t$, which further raises $\theta_t$, completing the loop. The loop
  is self-reinforcing when $|\Lambda| > 1$ (the slope of the fixed-point map
  exceeds one), which is guaranteed under the conditions in part (b).

  \textit{Assumptions used:} Assumption~\ref{ass:supply}(b) (endogenous HF
  supply), Assumption~\ref{ass:quality} ($\chi' > 0$),
  Assumption~\ref{ass:matching} (concave $M_t$), and the equilibrium price
  equation \eqref{eq:price-eq}. The multiple-equilibrium result additionally
  requires the local condition $|\Lambda(\hat\mu)| > 1$, which is a restriction
  on parameter magnitudes, not an additional assumption on the functional form.
  \hfill$\square$
\end{proof}

\begin{corollary}[Dollar Funding Run vs.\ Dollar Run]
  \label{cor:run-type}
  The run equilibrium in \cref{prop:run}(b) is characterized by high demand
  for conversion ($U_t$ large) and low quality of inside liquidity ($\chi_t$
  small), but not by a reduction in the aggregate quantity of dollar liquidity
  $L_t$. Agents continue to hold dollars --- both inside and outside --- but
  shift their demand toward safe outside dollars. This is a run \emph{within}
  the dollar ecosystem, not a run \emph{against} the dollar.
\end{corollary}

% ==============================================================================
\section{Link to the Empirical Strategy}
\label{sec:empirical-link}
% ==============================================================================

\cref{tab:theory-empirics} maps each theoretical proposition to its empirical
counterpart in the v8 design, the relevant prediction number, and the
assessment from the explorer feasibility report.

\begin{table}[htbp]
\centering
\begin{threeparttable}
\caption{Theory--Empirics Mapping}
\label{tab:theory-empirics}
\begin{tabular}{llll}
\toprule
Proposition & Mechanism & v8 Prediction & Feasibility \\
\midrule
\cref{prop:demand} & $\partial\mu^*/\partial U_t > 0$ & P1, P2 & Partial \\
\cref{prop:capacity} & $\partial\mu^*/\partial B_t < 0$ & P6 & Partial \\
\cref{prop:nonlinear} & Convexity in $\theta_t$ & P6 (heterogeneity) & Not feasible \\
\cref{prop:qe} & Quarter-end amplification & P7 & Partial \\
\cref{prop:swaplines} & Swap lines reduce spread & P6 (CF3) & Partial (event study) \\
\cref{prop:heterogeneity} & $|\beta_{FB}| < \beta_{HF}$ & P3, P4 & Partial \\
\cref{prop:run} & Multiple equilibria / run state & P6 (amplification) & Not directly testable \\
\bottomrule
\end{tabular}
\begin{tablenotes}\small
\item \textit{Notes:} Feasibility assessments are drawn from the explorer
feasibility report (score: 68/100, dated 2026-05-18). ``Partial'' means the
reduced-form prediction is testable but the structural channel cannot be
separately identified. ``Not feasible'' and ``not directly testable'' mean
the current data design does not support a clean identification of the
claimed mechanism; the paper should present these propositions as theoretical
benchmarks against which the data evidence is interpreted, not as claims that
are causally established.
\end{tablenotes}
\end{threeparttable}
\end{table}

% ==============================================================================
\section{Formal Extensions}
\label{sec:extensions}
% ==============================================================================

We sketch three extensions that would strengthen the theoretical contribution
for a top-5 journal submission. Each is stated precisely enough to be taken
to a formal proof with standard tools.

\subsection{Extension 1: Dynamics and Global Run Equilibrium}

\paragraph{Setup.} Make $U_t$ endogenous: agents observe the current $\chi_t$
and decide how urgently to seek conversion. Formally, let
\begin{equation}
  U_t = \bar{U} + \delta\,(1 - \chi_t) \cdot (U_t^{\max} - \bar{U}),
  \label{eq:Ut-endogenous}
\end{equation}
where $\delta \in [0,1]$ captures the sensitivity of conversion demand to
quality deterioration and $U_t^{\max}$ is the maximum mass that can seek
conversion. When $\delta = 0$, $U_t = \bar{U}$ (the static model). When
$\delta > 0$, quality deterioration feeds back into demand.

\paragraph{Main result.} Substituting \eqref{eq:Ut-endogenous} into the
equilibrium map and combining with $\chi_t = \chi(p_t)$ and $p_t =
\eta(U_t/B_t)^{-(1-\gamma)}$ yields a fixed-point equation in $\chi_t$
alone:
\[
  \chi_t = \chi\Bigl(\eta\Bigl[\bar{U} + \delta(1-\chi_t)(U_t^{\max}
  - \bar{U})\Bigr]^{-(1-\gamma)} B_t^{1-\gamma}\Bigr).
\]
Define the right-hand side as $\Gamma(\chi_t; B_t, \delta)$. A run equilibrium
exists alongside a normal equilibrium whenever $\Gamma$ has a slope
$|\partial\Gamma/\partial\chi| > 1$ at the normal equilibrium, which requires
$\delta$ and $(1-\gamma)$ large enough and $B_t$ small enough.

\paragraph{Tractability.} Fully tractable with standard fixed-point arguments
(Brouwer's theorem applied to the compact set $[\chi_{\min}, 1]$). The
comparative statics of the threshold in $\delta$ are computable in closed
form under $\chi(p) = p^{\alpha}$.

\subsection{Extension 2: Endogenous Intermediary Entry}

\paragraph{Setup.} Allow the mass of hedge funds to respond to $\mu_t$ via
free entry. Let the equilibrium HF mass $n_t^{HF}$ be determined by a
zero-profit condition: each HF earns $(\mu_t - c^{HF})\kappa/n_t^{HF}$ per
unit invested, and entry occurs until this equals the outside option $r > 0$.
This gives $n_t^{HF} = (\mu_t - c^{HF})\kappa/r$ when $\mu_t > c^{HF}$, and
zero otherwise. Total HF capacity: $B_t^{HF} = n_t^{HF} \cdot b^{HF} =
b^{HF}(\mu_t - c^{HF})\kappa/r$.

\paragraph{Main result.} Free entry flattens the HF supply curve: more entry
at higher $\mu_t$ provides additional stabilization. The condition for a
unique equilibrium under free entry is $\Lambda^{FE}(\mu) < 1$ where
$\Lambda^{FE}(\mu) = |\partial\mu^*/\partial B_t^{HF}|\cdot b^{HF}\kappa/r$.
Comparing to the fixed-mass case: free entry reduces $|\Lambda|$ whenever
$b^{HF}/r < \phi$ (marginal entrant supplies less capacity than the
representative incumbent), which need not hold when entry is cheap. The
model therefore predicts that endogenous entry can either dampen or amplify
the run, depending on the cost structure of HF participation.

\paragraph{Tractability.} Tractable under the demand-system specification:
the equilibrium price equation \eqref{eq:price-eq} is linear in $B_t^{HF}$,
and the free-entry condition is linear in $\mu_t$, so the system reduces to
two linear equations in $(\mu_t, B_t^{HF})$ with a closed-form solution.
Multiple equilibria do not arise under free entry if the system is everywhere
a contraction (sufficient condition: the free-entry slope $b^{HF}\kappa/r$
is small relative to $|B|$).

\subsection{Extension 3: Optimal Central Bank Intervention}

\paragraph{Setup.} Model the central bank as a player who chooses $B_t^{CB}
\in [0, \bar{B}^{CB}]$ at cost $c(B_t^{CB}) = c_0 (B_t^{CB})^2/2$ (convex
operational costs), maximizing social welfare $W = -\mu_t^* - \lambda
c(B_t^{CB})$, where $\lambda > 0$ is the weight on fiscal cost.

\paragraph{Main result.} The optimal intervention satisfies the first-order
condition:
\[
  -\frac{\partial\mu_t^*}{\partial B_t^{CB}} = \lambda c_0 B_t^{CB*}.
\]
By \cref{prop:capacity}, $|\partial\mu_t^*/\partial B_t^{CB}|$ is increasing
in $\theta_t$ (more strained markets benefit more from each unit of capacity).
Therefore the optimal $B_t^{CB*}$ is increasing in $\theta_t$: the central
bank should intervene more aggressively when the market is more stressed. This
is the theoretical foundation for the counter-cyclicality of swap lines
observed empirically.

Timing matters: intervening before the run equilibrium is reached (at $\mu_t
< \mu^{high}$) requires less capacity than restoring the normal equilibrium
after the run has started (at $\mu_t = \mu^{high}$). This generates an
\emph{option value of early intervention}: the welfare loss from waiting is
discrete (the system jumps to the run equilibrium) while the cost of early
intervention is continuous (gradual $B_t^{CB}$ provision). This delivers a
clear policy prediction: central banks should activate swap lines
pre-emptively, before the threshold in \cref{prop:run} is crossed.

\paragraph{Tractability.} Tractable under the linear demand system and with
the equilibrium price equation. The optimal intervention timing extends to a
dynamic problem under a finite-horizon stopping-time formulation, which is
standard in the financial crises literature (e.g., \citealt{DiasCosta2020}).
The stopping time is identified with the crossing of $\theta_t = \theta^*$
(the threshold in \cref{prop:nonlinear}).

% ==============================================================================
\section{What Is Proved vs.\ What Is Assumed}
\label{sec:summary}
% ==============================================================================

For transparency, we distinguish what is derived from model primitives (results)
from what is imposed on the model (assumptions).

\medskip

\noindent\textbf{Results (derived):}
\begin{itemize}[noitemsep]
  \item The signs $\partial\mu^*/\partial U_t > 0$ and $\partial\mu^*/\partial
        B_t < 0$ follow from the matching function and the CES cost-minimization
        problem, given $\chi'(p) > 0$.
  \item The convexity of $\mu^*$ in $\theta_t$ follows from the curvature of
        $p(\theta)$ and the strict monotonicity of $\chi(p)$, given $\chi'' > 0$.
  \item The quarter-end amplification follows from the parametric structure of
        FBO capacity in Assumption~\ref{ass:supply}(a).
  \item Multiple equilibria in \cref{prop:run} follow from the fixed-point
        structure of endogenous HF supply, given $|\Lambda| > 1$ locally.
\end{itemize}

\noindent\textbf{Assumptions (imposed):}
\begin{itemize}[noitemsep]
  \item $\chi'(p) > 0$ (Assumption~\ref{ass:quality}): quality is increasing in
        matching probability. This is a maintained assumption that defines the
        quality concept; it is consistent with but not derived from the model.
  \item $\chi'' > 0$ and $\Phi'' > 0$ (jointly): quality is convex in matching
        probability, and the spread is convex in quality. Both are additional
        conditions used only in \cref{prop:nonlinear}; both are stated
        explicitly there and neither is inherited from the earlier assumptions.
        Failure of either produces concavity or linearity, eliminating the
        threshold effect (see \cref{rem:convexity-conditions}).
  \item The CES specification \eqref{eq:CES} with $\sigma > 1$: the functional
        form of the aggregator. Relaxing this to a more general aggregator
        would preserve the qualitative results as long as inside and outside
        liquidity are gross substitutes.
  \item The linear demand system \eqref{eq:net-demand}: a local approximation.
        The global demand-system results (equilibrium price equation,
        counterfactuals) depend on this linearity; the matching-function results
        do not.
  \item $|\Lambda| > 1$ locally: the condition for multiple equilibria in
        \cref{prop:run}. This is a parameter restriction, not a functional-form
        assumption.
  \item The price-sensitivity ordering $|\beta_{FB}| < \beta_{HF}$: a
        \emph{maintained empirical hypothesis}, not a derived result (see
        \cref{rem:beta-ordering}). The $\beta_s$ parameters are free in the
        demand system; the ordering must be verified in estimation.
\end{itemize}

\noindent\textbf{What remains open:}
\begin{itemize}[noitemsep]
  \item A formal characterization of the basin of attraction of each equilibrium
        in the dynamic model (Extension 1). The conditions under which the
        economy converges to the run vs.\ normal equilibrium following a shock
        are not characterized here.
  \item A welfare analysis comparing the surplus lost in the run equilibrium
        to the fiscal cost of preemptive swap-line activation (Extension 3).
  \item An empirical test of the threshold effect in \cref{prop:nonlinear}: this
        requires a continuous $\hat\theta_t$ series, which is not available
        with current data (\emph{explorer feasibility report, Gap 3}).
\end{itemize}

% ==============================================================================
% References (placeholder -- to be populated from Bibliography_base.bib)
% ==============================================================================
\clearpage
\small
\begin{thebibliography}{99}

\bibitem[Bolton, Santos, and Scheinkman(2011)]{BoltonSantosScheinkman2011}
Bolton, P., T.\ Santos, and J.\ Scheinkman (2011).
\newblock Outside and inside liquidity.
\newblock \textit{Quarterly Journal of Economics}, 126(1), 259--321.

\bibitem[Goldsmith-Pinkham, Sorkin, and Swift(2020)]{GoldsmithPinkham2020}
Goldsmith-Pinkham, P., I.\ Sorkin, and H.\ Swift (2020).
\newblock Bartik instruments: What, when, why, and how.
\newblock \textit{American Economic Review}, 110(8), 2586--2624.

\bibitem[Hansen(2000)]{Hansen2000}
Hansen, B.\ E. (2000).
\newblock Sample splitting and threshold estimation.
\newblock \textit{Econometrica}, 68(3), 575--603.

\bibitem[Holmstr\"{o}m and Tirole(1998)]{HolmstromTirole1998}
Holmstr\"{o}m, B. and J.\ Tirole (1998).
\newblock Private and public supply of liquidity.
\newblock \textit{Journal of Political Economy}, 106(1), 1--40.

\bibitem[Koijen and Yogo(2019)]{KoijenYogo2019}
Koijen, R.\ S.\ J. and M.\ Yogo (2019).
\newblock A demand system approach to asset pricing.
\newblock \textit{Journal of Political Economy}, 127(4), 1475--1515.

\bibitem[Du, Tepper, and Verdelhan(2018)]{DuTepperVerdelhan2018}
Du, W., A.\ Tepper, and A.\ Verdelhan (2018).
\newblock Deviations from covered interest rate parity.
\newblock \textit{Journal of Finance}, 73(3), 915--957.

\bibitem[Rime, Schrimpf, and Syrstad(2022)]{RimeSchrimpfSyrstad2022}
Rime, D., A.\ Schrimpf, and O.\ Syrstad (2022).
\newblock Covered interest parity arbitrage.
\newblock \textit{Review of Financial Studies}, 35(11), 5185--5227.

\end{thebibliography}

\end{document}
